% you don't need a template
% use the following one

\documentclass[10pt]{article}
\usepackage[margin=2.5cm]{geometry}
\usepackage{booktabs}
\usepackage{float}

\title{Labwork 2\\\large HPC -- Get to know your GPU}
\author{
Le Nghiem Thanh Thuy - 2540023\\
University of Science and Technology of Hanoi
}
\date{}

\begin{document}
\maketitle
\section{Knowledge check}
I learned how to use CUDA for GPU computing in Python.
I used \texttt{is\_available()} to check whether CUDA is available.
I used \texttt{list\_devices()} to find available GPUs and \texttt{select\_device()} to choose one.
The Device provides the GPU’s name, ID, compute capability...
The Context provides information about the GPU’s memory...

\section{GPU information}

\begin{table}[H]
\centering
\caption{GPU information reported by \texttt{numba.cuda} on Google Colab.}
\label{tab:gpu}
\begin{tabular}{lr}
\toprule
\textbf{Property} & \textbf{Value} \\
\midrule
Device name & Tesla T4 \\
Compute capability & 7.5 \\
Multiprocessor count & 40 \\
Cores per SM & 64 \\
Total CUDA cores & 2560 \\
Total memory & 14.56 GB \\
\bottomrule
\end{tabular}
\end{table}


\end{document}